% Luc Destombe --- licence CC BY-NC-SA 4.0
% https://creativecommons.org/licenses/by-nc-sa/4.0/deed.fr
% Fichier autonome : compiler avec pdflatex, deux fois (signets, nombre de pages).
\def\lyceeversion{1.0}
\documentclass[a4paper,10pt]{article}
\makeatletter
\newif\iflycee@niveaux
\lycee@niveauxtrue
\newif\iflycee@cm
\newif\iflycee@palatino

\RequirePackage[utf8]{inputenc}
\RequirePackage[T1]{fontenc}
\RequirePackage[french]{babel}
\RequirePackage{etoolbox}

\newcommand{\ShorthandsOff}{\@ifundefined{shorthandoff}{}{\shorthandoff{;:!?}}}
\newcommand{\ShorthandsOn}{\@ifundefined{shorthandon}{}{\shorthandon{;:!?}}}
\ShorthandsOff
\AtBeginDocument{\ShorthandsOn}
\AtBeginEnvironment{tikzpicture}{\ShorthandsOff}

\RequirePackage{geometry}
\RequirePackage{fancyhdr}
\RequirePackage{lastpage}
\RequirePackage{multicol}
\RequirePackage{array}
\RequirePackage{enumitem}
\RequirePackage{xcolor}
\RequirePackage{needspace}

\RequirePackage{amsmath,amssymb}
\iflycee@cm\else\RequirePackage{fourier}\fi
\DeclareMathAlphabet{\mathcal}{OMS}{cmsy}{m}{n}
\renewcommand{\ttdefault}{lmtt}
\begingroup\catcode`\_=\active \catcode`\^=\active
\global\let_\sb \global\let^\sp
\endgroup
\AtBeginDocument{\catcode`\_=12 \mathcode`\_="8000
  \catcode`\^=12 \mathcode`\^="8000 }
\newlength{\retraitcalculs}
\setlength{\retraitcalculs}{2em}

\RequirePackage{tikz}
\usetikzlibrary{arrows.meta,calc,babel,shapes.misc}
\RequirePackage[most]{tcolorbox}
\tcbuselibrary{skins,breakable}

\definecolor{coulDef}{HTML}{1F4E79}
\definecolor{coulProp}{HTML}{7030A0}
\definecolor{coulRem}{HTML}{C55A11}
\definecolor{coulEx}{HTML}{2E7D32}
\definecolor{coulExo}{HTML}{B03A2E}
\definecolor{coulPart}{HTML}{1F4E79}
\definecolor{coulMeth}{HTML}{1B6E4F}
\definecolor{coulCourbe}{HTML}{0B5ED7}
\definecolor{coulCourbeB}{HTML}{C00000}
\definecolor{coulCourbeC}{HTML}{2E7D32}
\definecolor{coulGrille}{HTML}{8C8C8C}
\definecolor{coulRep}{HTML}{1B6E4F}
\definecolor{coulBar}{HTML}{2E7D32}

\newcommand{\R}{\mathbb{R}}
\newcommand{\Rs}{\mathbb{R}^{*}}
\newcommand{\N}{\mathbb{N}}
\newcommand{\Ns}{\mathbb{N}^{*}}
\newcommand{\Z}{\mathbb{Z}}
\newcommand{\Q}{\mathbb{Q}}
\renewcommand{\le}{\leqslant}
\renewcommand{\ge}{\geqslant}
\renewcommand{\approx}{\simeq}
\newcommand{\vect}[1]{\overrightarrow{#1}}
\newcommand{\norme}[1]{\left\lVert #1 \right\rVert}
\newcommand{\intervalle}[2]{\left[#1\,;#2\right]}
\RequirePackage{cancel}
\renewcommand{\CancelColor}{\color{coulExo}}
\newcommand{\fois}[1]{\times{\color{coulExo}#1}}
\newcommand{\barre}[1]{\cancel{#1}}
\newlength{\lycee@eqh}\newlength{\lycee@eqd}
\newsavebox{\lycee@eqa}\newsavebox{\lycee@eqb}\newsavebox{\lycee@eqc}
\newcommand{\lycee@eqcol}[1]{\hbox{$\renewcommand{\arraystretch}{1.35}%
  \begin{array}[t]{@{}>{\displaystyle}r@{}>{\displaystyle{}}l@{}}#1\end{array}$}}
\newcommand{\lycee@eqmesure}[1]{%
  \ifdim\ht#1>\lycee@eqh\setlength{\lycee@eqh}{\ht#1}\fi
  \ifdim\dp#1>\lycee@eqd\setlength{\lycee@eqd}{\dp#1}\fi}
\newcommand{\lycee@eqtrait}{\hspace{1.2em}ou\hspace{1.2em}}
\newcommand{\lycee@eqiff}{$\iff$\hspace{0.8em}}
\newcommand{\eqproduit}[2]{\par\smallskip\noindent C'est une équation produit :
  $A\times B=0 \iff A=0$ ou $B=0$.\par\nopagebreak\smallskip
  \begingroup\setlength{\lycee@eqh}{0pt}\setlength{\lycee@eqd}{0pt}%
  \setbox\lycee@eqa=\lycee@eqcol{#1}\setbox\lycee@eqb=\lycee@eqcol{#2}%
  \lycee@eqmesure\lycee@eqa\lycee@eqmesure\lycee@eqb%
  \noindent\hspace*{\retraitcalculs}\lycee@eqiff\usebox\lycee@eqa\lycee@eqtrait\usebox\lycee@eqb\endgroup\par\smallskip}
\newcommand{\eqproduittrois}[3]{\par\smallskip\noindent C'est une équation produit :
  $A\times B\times C=0 \iff A=0$ ou $B=0$ ou $C=0$.\par\nopagebreak\smallskip
  \begingroup\setlength{\lycee@eqh}{0pt}\setlength{\lycee@eqd}{0pt}%
  \setbox\lycee@eqa=\lycee@eqcol{#1}\setbox\lycee@eqb=\lycee@eqcol{#2}\setbox\lycee@eqc=\lycee@eqcol{#3}%
  \lycee@eqmesure\lycee@eqa\lycee@eqmesure\lycee@eqb\lycee@eqmesure\lycee@eqc%
  \noindent\hspace*{\retraitcalculs}\lycee@eqiff\usebox\lycee@eqa\lycee@eqtrait\usebox\lycee@eqb\lycee@eqtrait\usebox\lycee@eqc\endgroup
  \par\smallskip}

\tikzset{grille/.style={coulGrille,line width=0.4pt}}
\newcommand{\quadrillage}[4]{\draw[grille,step=1] (#1,#2) grid (#3,#4);}
\tikzset{
  croix/.style={cross out,draw,line width=0.6pt,inner sep=0pt,minimum size=4pt},
  croixdroite/.style={croix,rotate=45,line width=0.8pt},
}
\newcommand{\pt}[3]{\path (#1) node[croix]{} node[#2,font=\small,inner sep=2.5pt]{$#3$};}

\newcommand{\lycee@repere}[4]{%
  \draw[grille,step=1] (#1,#2) grid (#3,#4);
  \draw[-{Stealth[length=2mm]},thick] (#1,0)--({#3+0.4},0) node[below left=-1pt]{$x$};
  \draw[-{Stealth[length=2mm]},thick] (0,#2)--(0,{#4+0.4}) node[below left=-1pt]{$y$};
  \node[below left,font=\scriptsize,inner sep=1.5pt] at (0,0) {$0$};}
\newcommand{\lycee@gradx}[1]{\foreach \k/\l in {#1}{%
  \draw (\k,0.07)--(\k,-0.07) node[below,font=\scriptsize,inner sep=1.5pt]{$\l$};}}
\newcommand{\lycee@grady}[1]{\foreach \k/\l in {#1}{%
  \draw (0.07,\k)--(-0.07,\k) node[left,font=\scriptsize,inner sep=1.5pt]{$\l$};}}
\newcommand{\lycee@intff}[2]{\left[#1\,;#2\right]}
\newcommand{\lycee@intfo}[2]{\left[#1\,;#2\right[}
\newcommand{\lycee@intof}[2]{\left]#1\,;#2\right]}
\newcommand{\lycee@intoo}[2]{\left]#1\,;#2\right[}
\newcommand{\lycee@ens}[1]{\left\{#1\right\}}
\AtBeginDocument{%
  \providecommand{\repere}{\lycee@repere}%
  \providecommand{\gradx}{\lycee@gradx}%
  \providecommand{\grady}{\lycee@grady}%
  \providecommand{\Cf}{\mathcal{C}\sb{f}}%
  \providecommand{\Cg}{\mathcal{C}\sb{g}}%
  \providecommand{\intff}{\lycee@intff}%
  \providecommand{\intfo}{\lycee@intfo}%
  \providecommand{\intof}{\lycee@intof}%
  \providecommand{\intoo}{\lycee@intoo}%
  \providecommand{\ens}{\lycee@ens}}

\newlist{questions}{enumerate}{3}
\setlist[questions,1]{label=\arabic*\textdegree),leftmargin=*,itemsep=2pt,topsep=3pt}
\setlist[questions,2]{label=\alph*),leftmargin=*,itemsep=1pt,topsep=2pt}
\setlist[questions,3]{label=\roman*),leftmargin=*,itemsep=1pt,topsep=1pt}
\setlist[itemize]{label=\textbullet,leftmargin=*,itemsep=2pt,topsep=3pt}

\newcommand{\besoinplace}[1]{\par\penalty\@M\begingroup
  \dimen@=#1\relax\dimen@ii=\pagegoal\advance\dimen@ii-\pagetotal
  \ifdim\dimen@>\dimen@ii\ifdim\dimen@ii>\z@\vfil\fi\break\fi\endgroup}

\newcommand{\lycee@pied}{\thepage/\pageref{LastPage}}
\newcommand{\entete}[2]{%
  \pagestyle{fancy}\fancyhf{}%
  \fancyhead[L]{\footnotesize\color{coulPart}\textbf{#1}}%
  \fancyhead[R]{\footnotesize\color{coulPart}\textbf{#2}}%
  \fancyfoot[C]{\footnotesize\lycee@pied}%
  \renewcommand{\headrulewidth}{0.4pt}%
  \renewcommand{\headrule}{\hbox to\headwidth{%
    \color{coulPart!40}\leaders\hrule height \headrulewidth\hfill}}}

\RequirePackage{ccicons}
\newcommand{\lycee@licence}{{\scriptsize\color{gray}%
  \copyright\ Luc Destombe --- \ccbyncsa\ CC BY-NC-SA 4.0}}
\newif\iflycee@licence \lycee@licencetrue
\newcommand{\sanslicence}{\lycee@licencefalse}
\AtBeginDocument{\iflycee@licence\AddToHookNext{shipout/foreground}{%
  \put(0,0){\rlap{\kern\dimexpr1in+\hoffset+\oddsidemargin+\textwidth\relax
    \llap{\raisebox{-\dimexpr1in+\voffset+\topmargin+\headheight+\headsep
      +\textheight+\footskip\relax}{\lycee@licence}}}}}\fi}

\newcommand{\lycee@site}{faire-des-maths.fr}
\newcommand{\lycee@version}{{\scriptsize\color{gray}\lycee@site
  \ifdefined\lyceeversion\ --- Version~\lyceeversion\fi}}
\AtBeginDocument{\iflycee@licence\AddToHookNext{shipout/foreground}{%
  \put(0,0){\rlap{\kern\dimexpr1in+\hoffset+\oddsidemargin\relax
    \raisebox{-\dimexpr1in+\voffset+\topmargin+\headheight+\headsep
      +\textheight+\footskip\relax}{\lycee@version}}}}\fi}

\newcommand{\formulesserrees}{%
  \setlength{\abovedisplayskip}{3pt}\setlength{\belowdisplayskip}{3pt}%
  \setlength{\abovedisplayshortskip}{2pt}\setlength{\belowdisplayshortskip}{3pt}}

\setlength{\parindent}{0pt}

\RequirePackage[implicit=false,hidelinks,pdfencoding=auto,psdextra,bookmarksopen,
  bookmarksopenlevel=1,pdfcreator={},pdfproducer={}]{hyperref}
\RequirePackage{bookmark}
\hypersetup{pdfauthor={Luc Destombe},pdfkeywords={CC BY-NC-SA 4.0}}
\newcounter{lycee@signet}
\newcommand{\signet}[3][]{\stepcounter{lycee@signet}%
  \edef\lycee@dest{\if\relax\detokenize{#1}\relax
    signet.\arabic{lycee@signet}\else#1\fi}%
  \hypertarget{\lycee@dest}{}\bookmark[level=#2,dest=\lycee@dest]{#3}}
\pdfstringdefDisableCommands{%
  \def\R{R}\def\N{N}\def\Z{Z}\def\Q{Q}\def\vect#1{#1}\def\Cf{Cf}\def\Cg{Cg}\def\fois#1{×#1}%
  \def\barre#1{#1}\def\intervalle#1#2{[#1;#2]}\def\intff#1#2{[#1;#2]}%
  \def\textsuperscript#1{#1}\def\dfrac#1#2{#1/#2}\def\frac#1#2{#1/#2}%
  \def\vec#1{#1⃗}}%
\RequirePackage[official]{eurosym}

\geometry{top=0.8cm,bottom=1.3cm,left=1cm,right=1cm,footskip=0.6cm}

\pagestyle{fancy}
\fancyhf{}
\renewcommand{\headrulewidth}{0pt}
\fancyfoot[C]{\footnotesize Page \thepage{} / \pageref{LastPage}}

\newcommand{\titrefiche}[2]{%
  \begin{center}
    \Large\textbf{#1}\\[2pt]
    \large\textbf{#2}\\[2pt]
  \end{center}
  \vspace{0.10cm}}

\setlength{\columnseprule}{0.4pt}
\setlength{\columnsep}{15pt}
\renewcommand{\columnseprulecolor}{\color{black!45}}
\setlength{\multicolsep}{2pt}

\newcounter{partiefiche}
\renewcommand{\thepartiefiche}{\Roman{partiefiche}}
\newif\ifapresbandeau
\newcommand{\lycee@bandeau}[1]{%
  \par\penalty-600
  \vspace{0.08cm}%
  \noindent\colorbox{coulPart}{%
    \parbox{\dimexpr\linewidth-2\fboxsep\relax}{%
      \color{white}\bfseries\raggedright\strut#1\strut}}%
  \nopagebreak\par\nopagebreak
  \vspace{0.06cm}\nopagebreak
  \apresbandeautrue}
\newcommand{\partiefiche}[1]{%
  \refstepcounter{partiefiche}%
  \lycee@bandeau{\signet{1}{\thepartiefiche{} --- #1}\thepartiefiche{} --- #1}}
\newcommand{\partiefichesansnum}[1]{\lycee@bandeau{\signet{1}{#1}#1}}

\newcounter{exo}
\newlength{\espaceexo}\setlength{\espaceexo}{7pt}
\newlength{\espacetitre}\setlength{\espacetitre}{2pt}
\newcommand{\exo}[1][\relax]{%
  \par
  \ifapresbandeau\apresbandeaufalse\else\penalty-150\addvspace{\espaceexo}\fi
  \refstepcounter{exo}%
  \noindent\signet[exercice-\arabic{exo}]{2}{Exercice \arabic{exo}}%
  \textbf{\color{coulExo}\underline{Exercice \theexo}}%
  \ifx\relax#1\relax\else\textbf{ : #1}\fi
  \par\nopagebreak\vspace{\espacetitre}\nopagebreak
  \ignorespaces}

\setlist[enumerate]{nosep,leftmargin=*,itemsep=2pt,topsep=2pt}
\setlist[itemize]{nosep,leftmargin=*,topsep=2pt}
\newcommand{\choix}[3]{\par\hspace*{1em}%
  \makebox[0.3\linewidth][l]{a) #1}\makebox[0.3\linewidth][l]{b) #2}c) #3}
\newcommand{\deux}[2]{\par\hspace*{0.6em}%
  \makebox[0.48\linewidth][l]{#1}#2}
\newcommand{\trois}[3]{\par\hspace*{0.6em}%
  \makebox[0.32\linewidth][l]{#1}\makebox[0.32\linewidth][l]{#2}#3}

  \renewcommand{\exo}[1][\relax]{%
    \ifapresbandeau\apresbandeaufalse\else\par\penalty-150\fi
    \refstepcounter{exo}%
    \vspace{0.05cm}%
    \noindent\signet[exercice-\arabic{exo}]{2}{Exercice \arabic{exo}}%
    \textbf{\color{coulExo}\underline{Exercice \theexo}}%
    \ifx\relax#1\relax\else\textbf{ : #1}\fi%
    \par\nopagebreak
    \ignorespaces}
  \newcommand{\rappel}[1]{%
    \par\vspace{0.06cm}\nopagebreak
    \noindent\fcolorbox{black!35}{black!5}{%
      \parbox{\dimexpr\linewidth-2\fboxsep-2\fboxrule\relax}{%
        \small\textbf{Rappel.} #1}}%
    \par\vspace{0.06cm}\nopagebreak}
  \setlist[enumerate]{nosep,leftmargin=*}
  \setlist[itemize]{nosep,leftmargin=*}
  \setlength{\multicolsep}{3pt plus 1pt minus 1pt}
  \setlength{\premulticols}{10pt}
  \setlength{\postmulticols}{10pt}
  \newlist{expr}{enumerate}{1}
  \setlist[expr]{label={\Alph*\ $=$},nosep,leftmargin=*,itemsep=0pt}

\renewenvironment{center}{\par\addvspace{3pt}\centering}{\par\addvspace{3pt}}
\formulesserrees
\AtBeginDocument{\formulesserrees}
\clubpenalty=10000
\widowpenalty=10000
\displaywidowpenalty=10000
\brokenpenalty=10000
\setlength{\parskip}{0pt}

\RequirePackage{ragged2e}
\setlength{\RaggedRightRightskip}{0pt plus 1fil}
\AtBeginDocument{\RaggedRight\binoppenalty=10000 \relpenalty=10000 }
\g@addto@macro\@parboxrestore{\RaggedRight}
\makeatother
\usepackage{tkz-tab}
\geometry{top=0.8cm,bottom=1.3cm,left=1cm,right=1cm,footskip=0.7cm,headheight=12pt}

\newcommand{\e}{\mathrm{e}}

\begin{document}
\titrefiche{Terminale spécialité --- Limites de fonctions}{Fiche d'exercices du chapitre}

\begin{multicols}{2}

\partiefiche{Limite d'une fonction en l'infini}

\exo[Lire des limites et des asymptotes]
On donne les courbes $\Cf$ et $\Cg$ de deux fonctions $f$ et $g$ définies sur $\R$.
\begin{center}
\begin{tikzpicture}[x=0.32cm,y=0.32cm]
  \repere{-6}{-3}{6}{2}
  \gradx{-4,-2,2,4} \grady{-2,-1,1}
  \draw[coulCourbe,line width=1.1pt,domain=-6:6,samples=80] plot (\x,{-2+3/(1+\x*\x)});
  \node[coulCourbe,font=\footnotesize] at (1.6,1.4) {$\Cf$};
\end{tikzpicture}\hfill
\begin{tikzpicture}[x=0.32cm,y=0.32cm]
  \repere{-6}{-1}{3}{6}
  \gradx{-4,-2,2} \grady{1,2,3,4,5}
  \draw[coulCourbe,line width=1.1pt,domain=-6:1.8,samples=80] plot (\x,{2+0.5*exp(\x)});
  \node[coulCourbe,font=\footnotesize] at (-0.4,5.2) {$\Cg$};
\end{tikzpicture}
\end{center}
\begin{enumerate}
  \item Conjecturer les limites de $f$ et de $g$ en $-\infty$ et en $+\infty$.
  \item Pour chaque courbe, donner l'équation des asymptotes horizontales qu'elle semble
        admettre, en précisant en $-\infty$ ou en $+\infty$.
  \item Écrire en une phrase complète ce que signifie, pour la courbe $\Cg$, la limite de
        $g$ en $-\infty$.
\end{enumerate}

\exo[Conjecturer avec des valeurs]
Pour chaque fonction $f$, $g$, $h$, $k$ définie ci-dessous, calculer l'image de $10$, de
$100$ et de $1\,000$, puis de $-10$, de $-100$ et de $-1\,000$, et conjecturer les
limites en $+\infty$ et en $-\infty$. Interpréter graphiquement les limites finies.
\begin{enumerate}[label=\alph*)]
  \item $f$ définie pour tout nombre réel $x\neq 0$ par $f(x)=4-\dfrac{1}{x}$.
  \item $g$ définie sur $\R$ par $g(x)=x^{2}-3$.
  \item $h$ définie sur $\R$ par $h(x)=\dfrac{5}{x^{2}+1}$.
  \item $k$ définie sur $\R$ par $k(x)=1-x^{3}$.
\end{enumerate}

\exo[Un tableau de variations]
Voici le tableau de variations d'une fonction $f$ définie sur $\R$.
\begin{center}
\begin{tikzpicture}
  \tkzTabInit[lgt=1.4,espcl=2.2]{$x$/0.7,$f(x)$/1.4}{$-\infty$,$0$,$+\infty$}
  \tkzTabVar{-/$-1$,+/$3$,-/$2$}
\end{tikzpicture}
\end{center}
\begin{enumerate}
  \item Lire les limites de $f$ en $-\infty$ et en $+\infty$.
  \item En déduire, en phrases complètes, les asymptotes de la courbe $\Cf$.
  \item Tracer à main levée une courbe qui peut représenter $f$, avec ses asymptotes.
\end{enumerate}

\partiefiche{Limite d'une fonction en un nombre réel}

\exo[Conjecturer une limite en un nombre réel]
\begin{enumerate}
  \item Soit la fonction $f$ définie pour tout nombre réel $x\neq 1$ par
        $f(x)=\dfrac{4}{x-1}$.
  \begin{enumerate}[label=\alph*)]
    \item Calculer $f(1{,}1)$ ; $f(1{,}01)$ ; $f(1{,}001)$, puis $f(0{,}9)$ ;
          $f(0{,}99)$ ; $f(0{,}999)$.
    \item Conjecturer la limite à droite et la limite à gauche de $f$ en $1$, et en
          donner une interprétation graphique.
  \end{enumerate}
  \item Mêmes questions pour la fonction $g$ définie pour tout nombre réel $x\neq -2$ par
        $g(x)=\dfrac{1}{(x+2)^{2}}$, avec $x=-1{,}9$ ; $-1{,}99$ ; $-1{,}999$, puis
        $x=-2{,}1$ ; $-2{,}01$ ; $-2{,}001$.
\end{enumerate}

\exo[Lire deux asymptotes verticales]
La courbe $\Cf$ ci-dessous représente une fonction $f$ définie sur
$\R\setminus\ens{-1\,;2}$.
\begin{center}
\begin{tikzpicture}[x=0.4cm,y=0.4cm]
  \repere{-5}{-4}{6}{4}
  \gradx{-4,-3,-2,-1,1,2,3,4,5} \grady{-3,-2,-1,1,2,3}
  \begin{scope}
  \clip (-5,-4) rectangle (6,4);
  \draw[coulCourbe,line width=1.1pt,domain=-5:-1.07,samples=80] plot (\x,{1/((\x+1)*(\x-2))});
  \draw[coulCourbe,line width=1.1pt,domain=-0.92:1.92,samples=100] plot (\x,{1/((\x+1)*(\x-2))});
  \draw[coulCourbe,line width=1.1pt,domain=2.07:6,samples=80] plot (\x,{1/((\x+1)*(\x-2))});
  \end{scope}
  \node[coulCourbe,font=\footnotesize] at (3.6,1.4) {$\Cf$};
\end{tikzpicture}
\end{center}
\begin{enumerate}
  \item Lire les limites de $f$ en $-\infty$ et en $+\infty$.
  \item Lire la limite à gauche et la limite à droite de $f$ en $-1$, puis en $2$.
  \item Donner les équations des asymptotes à $\Cf$ que l'on peut conjecturer.
\end{enumerate}

\exo[Une valeur interdite dans un tableau]
Voici le tableau de variations d'une fonction $f$ définie sur
$\intoo{-\infty}{0}\cup\intoo{0}{+\infty}$.
\begin{center}
\begin{tikzpicture}
  \tkzTabInit[lgt=1.4,espcl=2.2]{$x$/0.7,$f(x)$/1.4}{$-\infty$,$0$,$+\infty$}
  \tkzTabVar{+/$1$,-D+/$-\infty$/$+\infty$,-/$1$}
\end{tikzpicture}
\end{center}
\begin{enumerate}
  \item Lire les limites de $f$ en $-\infty$, en $+\infty$, à gauche et à droite en $0$.
  \item En déduire, en phrases complètes, les asymptotes de la courbe $\Cf$.
  \item Tracer à main levée une courbe qui peut représenter $f$, avec ses asymptotes.
\end{enumerate}

\partiefiche{Opérations sur les limites}

\exo[Calculer par opérations en l'infini]
Pour chaque fonction, déterminer la limite demandée en citant la règle utilisée.
\begin{enumerate}[label=\alph*)]
  \item $f(x)=x^{3}+4x-1$, définie sur $\R$, en $+\infty$.
  \item $g(x)=\left(2-\dfrac{1}{x}\right)\left(x^{2}+3\right)$, définie pour $x\neq 0$,
        en $-\infty$.
  \item $h(x)=3-\dfrac{5}{x^{2}}$, définie pour $x\neq 0$, en $-\infty$.
  \item $k(x)=x\sqrt{x}+2$, définie sur $\intfo{0}{+\infty}$, en $+\infty$.
  \item $m(x)=\dfrac{-4}{x^{2}+1}$, définie sur $\R$, en $-\infty$.
  \item $p(x)=\dfrac{1}{x}-\sqrt{x}$, définie sur $\intoo{0}{+\infty}$, en $+\infty$.
\end{enumerate}

\exo[À gauche et à droite]
Pour chaque fonction, déterminer la limite à gauche et la limite à droite en la valeur
interdite.
\begin{enumerate}[label=\alph*)]
  \item $f(x)=\dfrac{3x+1}{x-2}$, définie pour tout nombre réel $x\neq 2$.
  \item $g(x)=\dfrac{1-x}{x+4}$, définie pour tout nombre réel $x\neq -4$.
  \item $h(x)=\dfrac{5}{(x-3)^{2}}$, définie pour tout nombre réel $x\neq 3$.
  \item $k(x)=\dfrac{x+1}{x^{2}}$, définie pour tout nombre réel $x\neq 0$.
\end{enumerate}

\exo[Avec un tableau de signes]
Soit la fonction $f$ définie sur $\R\setminus\ens{-1\,;1}$ par
$f(x)=\dfrac{2x+3}{x^{2}-1}$.
\begin{enumerate}
  \item Dresser le tableau de signes de $x^{2}-1$ sur $\R$.
  \item Déterminer les limites de $f$ à gauche et à droite en $-1$, puis en $1$.
  \item Déterminer les limites de $f$ en $-\infty$ et en $+\infty$.
  \item Interpréter graphiquement tous ces résultats.
\end{enumerate}

\exo[Le coût moyen d'une pièce imprimée]
Un atelier imprime en 3D des pièces de rechange pour vélos. Pour $x$ pièces, le coût
moyen d'une pièce, en euro, est $C(x)=3+\dfrac{200}{x}$, pour $x>0$.
\begin{enumerate}
  \item Calculer $C(100)$ et $C(1\,000)$.
  \item Déterminer la limite de $C$ en $+\infty$ et l'interpréter dans le contexte.
  \item Déterminer la limite de $C$ à droite en $0$ et l'interpréter dans le contexte.
  \item En déduire les asymptotes de la courbe de $C$.
\end{enumerate}

\partiefiche{Lever une indétermination}

\exo[Factoriser par le terme de plus haut degré]
Pour chaque fonction, indiquer la forme indéterminée, puis déterminer la limite demandée.
\begin{enumerate}[label=\alph*)]
  \item $f(x)=x^{2}-6x+1$, définie sur $\R$, en $+\infty$.
  \item $g(x)=-2x^{4}+3x^{2}+x$, définie sur $\R$, en $-\infty$.
  \item $h(x)=\dfrac{3x^{2}-x}{x^{2}+4}$, définie sur $\R$, en $+\infty$.
  \item $k(x)=\dfrac{5x-2}{x^{2}+1}$, définie sur $\R$, en $-\infty$.
  \item $m(x)=\dfrac{x^{2}+2x}{2x-1}$, définie pour $x\neq\frac{1}{2}$, en $+\infty$.
  \item $p(x)=\dfrac{4x^{3}-x}{2x^{3}+5}$, définie sur $\intof{-\infty}{-2}$, en $-\infty$.
\end{enumerate}

\exo[Une forme « $\frac{0}{0}$ » en un nombre réel]
Pour chaque fonction, factoriser le numérateur (et le dénominateur), simplifier, puis
déterminer la limite demandée.
\begin{enumerate}[label=\alph*)]
  \item $f(x)=\dfrac{x^{2}-16}{x-4}$, définie pour $x\neq 4$, en $4$.
  \item $g(x)=\dfrac{x^{2}-5x+6}{x-2}$, définie pour $x\neq 2$, en $2$.
  \item $h(x)=\dfrac{x^{2}+x-6}{x+3}$, définie pour $x\neq -3$, en $-3$.
  \item $k(x)=\dfrac{x^{2}-1}{x^{2}+x}$, définie pour $x\neq 0$ et $x\neq -1$, en $-1$.
\end{enumerate}

\exo[Avec la quantité conjuguée]
Pour chaque fonction, multiplier par la quantité conjuguée, puis déterminer la limite
demandée.
\begin{enumerate}[label=\alph*)]
  \item $f(x)=\sqrt{x+4}-\sqrt{x}$, définie sur $\intfo{0}{+\infty}$, en $+\infty$.
  \item $g(x)=\sqrt{x^{2}+1}-x$, définie sur $\R$, en $+\infty$.
  \item $h(x)=\dfrac{\sqrt{x}-3}{x-9}$, définie sur $\intfo{0}{9}\cup\intoo{9}{+\infty}$,
        en $9$.
\end{enumerate}

\exo[Deux valeurs interdites]
Soit la fonction $f$ définie par $f(x)=\dfrac{2x^{2}-x-1}{x^{2}-1}$.
\begin{enumerate}
  \item Déterminer l'ensemble de définition de $f$.
  \item Vérifier que $2x^{2}-x-1=(x-1)(2x+1)$, puis simplifier $f(x)$ en précisant
        pour quels $x$ l'égalité obtenue est vraie.
  \item Déterminer les limites de $f$ en $-\infty$ et en $+\infty$.
  \item Déterminer la limite de $f$ en $1$, puis les limites à gauche et à droite en
        $-1$.
  \item La droite d'équation $x=1$ est-elle asymptote à $\Cf$ ? Et la droite d'équation
        $x=-1$ ? Justifier.
\end{enumerate}

\partiefiche{Asymptotes et position relative}

\exo[Asymptotes et position]
Soit la fonction $f$ définie pour tout nombre réel $x\neq -2$ par
$f(x)=\dfrac{3x-1}{x+2}$, et $\Cf$ sa courbe.
\begin{enumerate}
  \item Montrer que, pour tout nombre réel $x\neq -2$, $f(x)=3-\dfrac{7}{x+2}$.
  \item Déterminer les limites de $f$ en $-\infty$ et en $+\infty$, et en déduire une
        asymptote à $\Cf$.
  \item Déterminer les limites de $f$ à gauche et à droite en $-2$, et en déduire une
        asymptote à $\Cf$.
  \item Étudier la position de $\Cf$ par rapport à la droite d'équation $y=3$.
\end{enumerate}

\exo[Toujours au-dessus]
Soit la fonction $g$ définie sur $\R$ par $g(x)=\dfrac{2x^{2}+3}{x^{2}+1}$, et $\Cg$
sa courbe.
\begin{enumerate}
  \item Montrer que, pour tout nombre réel $x$, $g(x)=2+\dfrac{1}{x^{2}+1}$.
  \item Déterminer les limites de $g$ en $-\infty$ et en $+\infty$. Interpréter
        graphiquement.
  \item Étudier la position de $\Cg$ par rapport à son asymptote.
  \item Expliquer pourquoi $\Cg$ n'a pas d'asymptote verticale.
\end{enumerate}

\exo[La vitesse moyenne d'une cycliste]
Une cycliste monte un col de $d$~km à la vitesse moyenne de $15$~km/h, puis redescend
par la même route à la vitesse moyenne de $v$~km/h, avec $v>0$.
\begin{enumerate}
  \item Exprimer, en fonction de $d$ et de $v$, la durée de la montée, celle de la
        descente, puis montrer que sa vitesse moyenne sur tout le trajet est
        $V(v)=\dfrac{30v}{v+15}$.
  \item Calculer $V(15)$ et $V(45)$.
  \item Déterminer la limite de $V$ en $+\infty$.
  \item Montrer que, pour tout $v>0$, $30-V(v)=\dfrac{450}{v+15}$.
  \item En déduire que, quelle que soit sa vitesse dans la descente, la cycliste ne
        peut pas atteindre une moyenne de $30$~km/h. Interpréter graphiquement.
\end{enumerate}

\partiefiche{Limite d'une fonction composée}

\exo[Calculer la limite d'une composée]
Pour chaque fonction, déterminer la limite demandée en posant $X=\ldots$
\begin{enumerate}[label=\alph*)]
  \item $f(x)=\sqrt{4-\dfrac{1}{x}}$, définie sur $\intoo{1}{+\infty}$, en $+\infty$.
  \item $g(x)=\sqrt{x^{2}+x+1}$, définie sur $\R$, en $-\infty$.
  \item $h(x)=\dfrac{1}{\left(3+\frac{1}{x}\right)^{2}}$, définie sur
        $\intoo{-\infty}{-1}$, en $-\infty$.
  \item $k(x)=\sqrt{\dfrac{2x+1}{x-3}}$, définie sur $\intoo{3}{+\infty}$, en $+\infty$.
\end{enumerate}

\exo[Une racine au dénominateur]
Soit la fonction $f$ définie sur $\intoo{2}{+\infty}$ par $f(x)=\dfrac{1}{\sqrt{x-2}}$.
\begin{enumerate}
  \item Déterminer la limite de $f$ à droite en $2$.
  \item Déterminer la limite de $f$ en $+\infty$.
  \item En déduire les asymptotes de la courbe $\Cf$.
\end{enumerate}

\exo[La racine d'un quotient]
Soit la fonction $g$ définie pour tout nombre réel $x\neq -3$ par
$g(x)=\dfrac{2x-1}{x+3}$.
\begin{enumerate}
  \item Dresser le tableau de signes de $g(x)$, puis en déduire l'ensemble $D$ des
        nombres réels $x$ pour lesquels $g(x)\geq 0$.
  \item Déterminer les limites de $g$ en $-\infty$ et en $+\infty$.
  \item Déterminer la limite de $g$ à gauche en $-3$.
  \item Soit $f$ la fonction définie sur $D$ par $f(x)=\sqrt{\dfrac{2x-1}{x+3}}$.
        Déterminer les limites de $f$ en $-\infty$, en $+\infty$ et à gauche en $-3$.
  \item En déduire les asymptotes de la courbe de $f$.
\end{enumerate}

\exo[D'après bac 2026, Asie : vrai ou faux ?]
Indiquer si l'affirmation suivante est vraie ou fausse, en justifiant.

On considère la fonction $f$ définie sur l'intervalle $\intoo{1}{+\infty}$ par :
$f(x)=\dfrac{x-1}{\sqrt{x^{2}-1}}$.

\textbf{Affirmation :} la fonction $f$ admet pour limite $1$ en $+\infty$.

\partiefiche{Limites et comparaison}

\exo[Comparaison ou gendarmes]
Pour chaque fonction, déterminer la limite en $+\infty$ en justifiant par un théorème
de comparaison ou par le théorème des gendarmes.
\begin{enumerate}[label=\alph*)]
  \item $f(x)=x^{2}+3\sin(x)$, définie sur $\R$.
  \item $g(x)=-3x+2\cos(x)$, définie sur $\R$.
  \item $h(x)=\dfrac{\cos(x)}{x}$, définie sur $\intoo{0}{+\infty}$.
  \item $k(x)=5+\dfrac{\sin(x)}{x^{2}}$, définie sur $\intoo{0}{+\infty}$.
\end{enumerate}

\exo[En $+\infty$ et en $-\infty$]
Soit la fonction $f$ définie sur $\R$ par $f(x)=2x+3\cos(x)$.
\begin{enumerate}
  \item Montrer que, pour tout nombre réel $x$, $2x-3\leq f(x)\leq 2x+3$.
  \item En déduire la limite de $f$ en $+\infty$.
  \item En déduire la limite de $f$ en $-\infty$.
\end{enumerate}

\exo[Encadrer une racine carrée]
Soit la fonction $f$ définie sur $\intoo{0}{+\infty}$ par $f(x)=\dfrac{\sqrt{x^{2}+1}}{x}$.
\begin{enumerate}
  \item Montrer que, pour tout $x>0$, $x^{2}\leq x^{2}+1\leq (x+1)^{2}$.
  \item En déduire que, pour tout $x>0$, $1\leq f(x)\leq\dfrac{x+1}{x}$.
  \item Déterminer la limite de $\dfrac{x+1}{x}$ en $+\infty$, puis celle de $f$.
  \item Interpréter graphiquement le résultat.
\end{enumerate}

\exo[D'après bac 2026, Antilles-Guyane : vrai ou faux ?]
Pour chaque affirmation, indiquer si elle est vraie ou fausse, en justifiant.
\begin{enumerate}[label=\alph*)]
  \item Si une fonction $f$ vérifie $f(x)\geq x^{2}$ pour tout $x>5$, alors
        $\lim\limits_{x\to+\infty}f(x)=+\infty$.
  \item On considère la fonction $f$ définie sur l'intervalle $\intoo{0}{+\infty}$ par
        $f(x)=1+\dfrac{\sin(x)}{x}$. La droite d'équation $y=1$ est une asymptote
        horizontale de la courbe représentative de la fonction $f$.
  \item Une fonction qui a pour limite $+\infty$ en $+\infty$ est croissante sur un
        intervalle $\intoo{a}{+\infty}$. (On pourra comparer les images de
        $2k\pi+\frac{\pi}{2}$ et de $2k\pi+\pi$ par $x\mapsto x+2\sin(x)$, pour
        $k\in\N$.)
\end{enumerate}

\partiefiche{Limites de la fonction exponentielle}

\exo[Avec des exponentielles]
Pour chaque fonction, définie sur $\R$ (sur $\R\setminus\ens{0}$ pour $p$), déterminer
les limites demandées.
\begin{enumerate}[label=\alph*)]
  \item $f(x)=\e^{2x+1}$, en $-\infty$ et en $+\infty$.
  \item $g(x)=7-\e^{x}$, en $-\infty$ et en $+\infty$.
  \item $h(x)=\e^{-x^{2}}$, en $+\infty$.
  \item $k(x)=(3-x)\,\e^{-x}$, en $-\infty$.
  \item $m(x)=x\left(1+\e^{-x}\right)$, en $-\infty$ et en $+\infty$.
  \item $p(x)=\e^{\frac{1}{x}}$, en $+\infty$, puis à gauche et à droite en $0$.
\end{enumerate}

\exo[Factoriser par $\e^{x}$]
Soit la fonction $f$ définie sur $\R$ par $f(x)=\dfrac{2\e^{x}-3}{\e^{x}+4}$, et $\Cf$
sa courbe.
\begin{enumerate}
  \item Déterminer la limite de $f$ en $-\infty$.
  \item Montrer que, pour tout nombre réel $x$, $f(x)=\dfrac{2-3\e^{-x}}{1+4\e^{-x}}$.
  \item En déduire la limite de $f$ en $+\infty$.
  \item Interpréter graphiquement les deux limites.
\end{enumerate}

\exo[Avec le nombre $\e$]
\begin{enumerate}
  \item Soit la suite $(u_{n})$ définie pour tout $n\in\N$ par $u_{n}=\e^{n}$. Montrer
        que $(u_{n})$ est géométrique, puis déterminer sa limite.
  \item Soit la suite $(v_{n})$ définie pour tout $n\in\N$ par $v_{n}=\e^{-n}$.
        Déterminer sa limite de deux façons : comme suite géométrique, puis avec la
        limite de $\e^{x}$ en $-\infty$.
  \item Soit la suite $(w_{n})$ définie pour tout $n\in\N$ par
        $w_{n}=\dfrac{3\e^{n}-1}{\e^{n}+2}$. Déterminer sa limite.
\end{enumerate}

\exo[La charge d'une batterie]
On branche la batterie déchargée d'un vélo électrique. L'énergie stockée, en
wattheure (Wh), $t$ heures après le branchement, est modélisée par
$Q(t)=500\left(1-\e^{-0{,}4t}\right)$, pour $t\geq 0$.
\begin{enumerate}
  \item Calculer $Q(0)$, puis une valeur approchée de $Q(5)$ à l'unité.
  \item Déterminer la limite de $Q$ en $+\infty$ et l'interpréter dans le contexte.
  \item Calculer $Q'(t)$ et en déduire le sens de variation de $Q$ sur
        $\intfo{0}{+\infty}$.
  \item Étudier la position de la courbe de $Q$ par rapport à son asymptote.
        Interpréter dans le contexte.
\end{enumerate}

\partiefiche{Croissances comparées}

\exo[Calculer par croissances comparées]
Pour chaque fonction, définie sur $\R$ (sur $\intoo{0}{+\infty}$ pour $k$), déterminer
la limite demandée.
\begin{enumerate}[label=\alph*)]
  \item $f(x)=\e^{x}-x^{2}$, en $+\infty$.
  \item $g(x)=\left(x^{2}+1\right)\e^{-x}$, en $+\infty$.
  \item $h(x)=(x-3)\,\e^{x}$, en $-\infty$.
  \item $k(x)=\dfrac{\e^{x}}{x^{3}}$, en $+\infty$.
  \item $m(x)=x+2-\e^{x}$, en $+\infty$.
  \item $p(x)=\e^{2x}-x\e^{x}$, en $+\infty$.
\end{enumerate}

\exo[Croissances comparées et composée]
\begin{enumerate}
  \item Soit la fonction $f$ définie sur $\R$ par $f(x)=x\e^{-2x}$.
  \begin{enumerate}[label=\alph*)]
    \item Vérifier que, pour tout nombre réel $x$, $f(x)=\dfrac{1}{2}\times 2x\,\e^{-2x}$.
    \item En posant $X=2x$, déterminer la limite de $f$ en $+\infty$.
  \end{enumerate}
  \item Soit la fonction $g$ définie pour tout nombre réel $x\neq 0$ par
        $g(x)=\dfrac{\e^{2x}}{x}$. Déterminer les limites de $g$ en $-\infty$ et en
        $+\infty$.
\end{enumerate}

\exo[D'après bac 2026, Polynésie : vrai ou faux ?]
Pour chaque affirmation, indiquer si elle est vraie ou fausse, en justifiant.
\begin{enumerate}[label=\alph*)]
  \item Soit $f$ la fonction définie sur $\R$ par $f(x)=(x+1)\e^{-x}$. La limite de la
        fonction $f$ en $+\infty$ est $+\infty$.
  \item $\lim\limits_{x\to-\infty}x^{2}\e^{x}=+\infty$.
  \item $\lim\limits_{x\to+\infty}\left(\e^{x}-x^{10}\right)=+\infty$.
\end{enumerate}

\exo[Une étude complète]
Soit la fonction $f$ définie sur $\R$ par $f(x)=(x-1)\,\e^{x}$, et $\Cf$ sa courbe.
\begin{enumerate}
  \item Déterminer la limite de $f$ en $-\infty$. Interpréter graphiquement.
  \item Déterminer la limite de $f$ en $+\infty$.
  \item Montrer que, pour tout nombre réel $x$, $f'(x)=x\e^{x}$.
  \item Dresser le tableau de variations de $f$, limites comprises.
  \item Étudier la position de $\Cf$ par rapport à son asymptote.
\end{enumerate}

\partiefichesansnum{Exercices type bac}

\exo[D'après bac 2026, Métropole]
Un essai clinique est réalisé auprès de plusieurs patients volontaires afin de valider
ou non l'efficacité d'un médicament.

On demande au patient d'avaler le médicament sous la forme d'un comprimé de 1 mg puis
on étudie la quantité de ce médicament présente dans le sang au cours du temps.

La quantité de ce médicament en fonction du temps écoulé après la prise du comprimé, en
mg, est modélisée par la fonction $f$ dérivable sur $\intfo{0}{+\infty}$, définie par :
\[f(t)=4t\e^{-\frac{3}{2}t},\quad t\ \text{étant exprimé en heures}.\]
\begin{enumerate}
  \item Déterminer $\lim\limits_{t\to+\infty}f(t)$.
        Interpréter ce résultat dans le contexte de l'exercice.
  \item
  \begin{enumerate}[label=\alph*)]
    \item Vérifier en détaillant les calculs que pour tout réel $t$ positif,
          $f'(t)=(4-6t)\e^{-\frac{3}{2}t}$.
    \item Étudier le signe de la dérivée et en déduire le tableau de variations de la
          fonction $f$ sur $\intfo{0}{+\infty}$.
    \item Au bout de combien de temps la quantité de ce médicament présente dans le sang
          est-elle maximale ?
  \end{enumerate}
\end{enumerate}

\exo[D'après bac 2025, Amérique du Nord]
On considère la fonction $f$ définie pour tout nombre réel $x$ par :
\[f(x)=\left(x^{2}+3x+2\right)\e^{-x}.\]
\begin{enumerate}
  \item Démontrer que la limite de la fonction $f$ en $+\infty$ est égale à $0$.
  \item Démontrer que la limite de la fonction $f$ en $-\infty$ est égale à $+\infty$.
  \item On admet que la fonction $f$ est dérivable sur $\R$. On note $f'$ la fonction
        dérivée de $f$ sur $\R$.
  \begin{enumerate}[label=\alph*)]
    \item Vérifier que, pour tout nombre réel $x$, $f'(x)=\left(-x^{2}-x+1\right)\e^{-x}$.
    \item Déterminer le signe de la fonction dérivée $f'$ sur $\R$ puis en déduire les
          variations de la fonction $f$ sur $\R$.
  \end{enumerate}
  \item Expliquer pourquoi la fonction $f$ est positive sur l'intervalle
        $\intfo{0}{+\infty}$.
\end{enumerate}

\exo[D'après bac 2025, Centres étrangers]
\textbf{Partie A}

On considère la fonction $f$ définie sur l'intervalle $\intfo{0}{+\infty}$ par :
\[f(x)=\dfrac{1}{a+\e^{-bx}}\]
où $a$ et $b$ sont deux constantes réelles strictement positives.

On admet que la fonction $f$ est dérivable sur l'intervalle $\intfo{0}{+\infty}$.
La fonction $f$ admet pour représentation graphique la courbe $\Cf$ ci-dessous :
\begin{center}
\begin{tikzpicture}[x=0.2cm,y=3.5cm]
  \foreach \k in {0,5,...,30} \draw[grille] (\k,0)--(\k,1.1);
  \foreach \k in {0,0.1,...,1.05} \draw[grille] (0,\k)--(33,\k);
  \draw[-{Stealth[length=2mm]},thick] (0,0)--(33.5,0) node[below left=-1pt]{$x$};
  \draw[-{Stealth[length=2mm]},thick] (0,0)--(0,1.15) node[below left=-1pt]{$y$};
  \foreach \k in {5,10,...,30} \draw (\k,0.02)--(\k,-0.02) node[below,font=\scriptsize]{$\k$};
  \draw (0.4,0.5)--(-0.4,0.5) node[left,font=\scriptsize]{$0{,}5$};
  \draw (0.4,1)--(-0.4,1) node[left,font=\scriptsize]{$1$};
  \node[below left,font=\scriptsize,inner sep=1.5pt] at (0,0) {$0$};
  \draw[coulCourbeB,line width=0.7pt] (0,0.5)--(12,1.1);
  \draw[coulCourbe,line width=1.1pt,domain=0:33,samples=80] plot (\x,{1/(1+exp(-0.2*\x))});
  \path (0,0.5) node[croix]{} node[below right,font=\footnotesize]{A};
  \path (10,1) node[croix]{} node[below right,font=\footnotesize]{B};
  \node[coulCourbe,font=\footnotesize] at (29,0.9) {$\Cf$};
\end{tikzpicture}
\end{center}
On considère les points $A(0\,;0{,}5)$ et $B(10\,;1)$.

On admet que la droite $(AB)$ est tangente à la courbe $\Cf$ au point $A$.
\begin{enumerate}
  \item Par lecture graphique, donner une valeur approchée de $f(10)$.
  \item On admet que $\lim\limits_{x\to+\infty}f(x)=1$. Donner une interprétation
        graphique de ce résultat.
  \item Justifier que $a=1$.
  \item Déterminer le coefficient directeur de la droite $(AB)$.
  \item
  \begin{enumerate}[label=\alph*)]
    \item Déterminer l'expression de $f'(x)$ en fonction de $x$ et de la constante $b$.
    \item En déduire la valeur de $b$.
  \end{enumerate}
\end{enumerate}

\textbf{Partie B}

On admet, dans la suite de l'exercice, que la fonction $f$ est définie sur
l'intervalle $\intfo{0}{+\infty}$ par :
\[f(x)=\dfrac{1}{1+\e^{-0{,}2x}}\]
\begin{enumerate}
  \item Déterminer $\lim\limits_{x\to+\infty}f(x)$.
  \item Étudier les variations de la fonction $f$ sur l'intervalle $\intfo{0}{+\infty}$.
\end{enumerate}

\end{multicols}

\end{document}
